\section{Discussion}
In this section, we present a summary of the key findings per research questions, provide the implications for both research and practice, and address the limitations of our study. Additionally, we outline potential directions for future research.

\subsection{Summary of Results and Research Questions}
\subsubsection{RQ1: Effectiveness} Participants demonstrated higher \textit{effectiveness} scores on the DS-enhanced visualisations compared to the conventional visualisations, with a statistically significant difference and a small effect size. Interestingly, the \textit{effectiveness} scores were higher for the first three levels of Bloom's taxonomy—\textit{Identify}, \textit{Understand}, and \textit{Apply}—when using DS-enhanced visualisations. These results are consistent with those of \citet{shao2024datastorytelling}, who also observed that DS significantly improved participants' ability to identify and comprehend key data insights in the first two levels of cognition. However, unlike previous work \cite{shao2024datastorytelling, zdanovic2022influence}, we studied the effectiveness at further levels of Bloom's taxonomy. For tasks requiring higher-order cognitive skills, such as \textit{Analyse} and \textit{Evaluate}, the \textit{effectiveness} scores were higher for the conventional visualisations. Although the differences were not significant, this may suggest that DS elements are most effective in simpler tasks, where narrative elements help structure information, but less so when participants must engage in more abstract thinking, where the narrative may not provide additional value or could even distract from critical thinking processes. The more structured nature of DS-enhanced visualisations, which often highlight specific data points, might direct users' attention and thinking, whereas conventional visualisations offer more flexibility for individual interpretation. 

\subsubsection{RQ2: Efficiency} Regarding \textit{efficiency}, participants showed no significant overall difference between DS-enhanced and conventional visualisations. However, at higher cognitive levels (\textit{Understand}, \textit{Apply}, \textit{Analyse}, \textit{Evaluate}), participants were significantly more efficient with DS-enhanced visualisations while for conventional visualisations they were slightly quicker at the \textit{Identify} level. Our findings from the \textit{Identify} level are in line with \citet{shao2024datastorytelling} where at tasks of lower cognitive levels, like \textit{Identify}, the DS elements on visualisations made participants take longer to answer questions correctly. One possible explanation is that the additional narrative and contextual elements introduced in DS-enhanced visualisations can distract or delay participants when the task requires simple recall or recognition. In contrast, conventional visualisations, being more straightforward and minimalistic, allow for quicker retrieval of known information. 

Another possible explanation may be related to the hierarchical ordering of the tasks. All the participants were presented tasks following the order Bloom's taxonomy, meaning that the first question was an \textit{Identify} question followed by \textit{Understand} and higher cognitive levels. When participants were exposed to the DS-enhanced visualisations, they may have taken longer in the \textit{Identify} tasks than in the conventional visualisation condition as they were reviewing all of the information available in the visualisation. Once they had understood the story they were faster at answering correctly the questions from all the higher levels (\textit{Understand}, \textit{Apply}, \textit{Analyse}, \textit{Evaluate}). This is consistent with the discussion from \citet{liu2020survey} and \citet{galati2021} who suggest that individuals will take time to engage in sensemaking behaviours with visualisations in order to understand the contents before they complete any tasks using the visualisation's information. The findings from this study also resonates with the argument put forth by \citet{ryan2016visual} who believes that individuals will typically notice visual attributes, such as colour, before focusing on the deeper meaning of the visualisation. It is possible that participants took the time to perceive the artistic visual elements of the DS-enhancements before reviewing the intended message of the data story.

\subsubsection{RQ3: Quality of Complex Creation Tasks} Participants' text composition showed no significant overall difference between DS-enhanced and conventional visualisations. However, conventional visualisations prompted more sentences at the \textit{Identify} and \textit{Evaluate} levels, while DS-enhanced visualisations led to more sentences in the \textit{Understand}, \textit{Analyse}, and \textit{Create} categories. The pattern of engaging in different levels of cognition depending on the added elements on the visualisation resembles the findings from \citet{stokes2023striking}, where participants engaged in deeper semantic content takeaways when the text was placed near relevant data points on the chart, rather than on the title.  In our study, contextual information on data points in the DS-enhanced visualisations may have contributed to participants engaging more with higher-level cognitive tasks, providing further evidence of the relationship between text placement and cognitive depth in visualisation-based tasks.

One possible explanation for the higher number of \textit{Identify} sentences in the conventional visualisation condition is that participants may have felt the need to restate and emphasise information that was less visually highlighted, as opposed to DS-enhanced where key values were already emphasised. This could have led participants in the DS-enhanced condition to move beyond simple recall and engage more frequently at the \textit{Understand} level and reach higher cognitive levels, such as \textit{Apply}, \textit{Analyse}, and \textit{Create}. Interestingly, participants using conventional visualisations produced more \textit{Evaluate} sentences. Similar to the discussion on RQ1, this could be attributed to conventional visualisations being more open-ended, thus encouraging interpretation and critical thinking. In contrast, DS-enhanced visualisations, which are designed to guide users toward a specific message or narrative, may reduce the need for critical evaluation, such as identifying missing elements that could influence interpretation. This hypothesis aligns with the work by \citet{kong2018} on the effect of visualisation titles which found that participants were often unaware of the potential bias in visualisations leading to misconceptions of information. Our work extends their findings suggesting that data storytelling elements make the visualisation more convincing and there is less awareness of the limitations or biases of the information presented and the way it is presented. \citet{rogha2024impact} hypothesised that the passive data interpretation could be due to the cognitive load induced by the narrative elements.  In the work of \citet{garreton2024attitudinal}, visualisations are effective in eliciting emotional impact and driving attitude change. Building on this, we argue that while the persuasive power of such visualisations can be instrumental in driving action, such as addressing urgent issues like the water crisis \cite{garreton2024attitudinal}, this same persuasiveness may have unintended consequences. Specifically, these visualisations could inadvertently limit users’ ability to question the presented narrative, potentially leading to over-reliance on the guided interpretation provided.  


\subsubsection{RQ4: User Perception} Lastly, regarding user preferences, the majority of participants preferred DS-enhanced over conventional visualisations for the placement of contextual information and clear distinctions between countries, while conventional visualisations were preferred for their simplicity and minimal clutter. This aligns with the findings by \citet{yen2020} where participants perceived visualisations as useful when the text was presented in shorter segments that focused on one aspect at a time rather than all together. The idea that DS enhancements, such as annotations and explanatory titles, help to contextualise information was shared by the participants from \citet{echeverria2018exploratory, stokes2022text, milesi2024storytelling} who praised textual elements for their ability to make insights understandable and accessible. Similarly, \citet{stokes2023striking} found that participants preferred charts with a narrative or story to a text-only version of the data or visualisations without text annotations. A limitation of their work was the use of synthetically generated line charts, which lacked the complexity of real-world examples. Our study addresses this gap by providing empirical validation of \citet{stokes2022text,stokes2023striking} observations with real-world datasets, particularly emphasising the importance of placing contextual information near relevant data points. Additionally, the results are congruent with \citet{bateman2010} where visualisations with DS elements were perceived as more attractive, enjoyable, and easier to read and remember. Lastly, it is worth noting that there are individual preferences for information presentation or omission \cite{lundgard2022accessible}. \citet {lundgard2022accessible} advocate for adaptable descriptions that balance elemental, perceptual, and contextual information to meet diverse user needs in natural language descriptions of visualisations. We propose that this principle also applies to contextual annotations in DS-enhanced visualisations.

\subsection{Implications for Research}
In the literature, DS-enhancements have been a focal point for their perceived ability to facilitate the sensemaking process and more effectively and efficiently guide individuals towards complex insights compared to conventional data visualisations \cite{segel2010narrative, boy2015storytelling, Daradkeh2021}. This study, along with the work of \citet{shao2024datastorytelling}, has suggested that with respect to effectiveness, this idea holds for tasks that require a low level of cognitive processing. However, our work indicates that for tasks that require higher order thinking DS-enhancement may be less effective at facilitating the sensemaking process such that individuals can uncover patterns and insights on their own \cite{campos2021makingsense}. This indicates that the complexity of a task can significantly impact the \textit{effectiveness} of DS-enhanced visualisations. This presents an interesting direction for future works in establishing whether the nature of the task (e.g., written, verbal, etc.) can impact the effectiveness just as much as the complexity of the task.

In terms of efficiency, while \citet{shao2024datastorytelling} observed no improvements in speed between DS-enhanced and conventional visualisations, this study indicated that DS-enhancements can significantly impact how quickly individuals are able to extract insights from complex data. This is consistent with the findings from \citet{echeverria2018exploratory, de2014evaluating}, and \citet{Pozdniakov2023how} who argue that DS-enhancements enable a deeper understanding of the presented narrative and thus facilitate faster interpretation of the material. This research could be expanded to investigate the specific DS-enhancements that aid in efficient extraction of insights or examine whether the improvements in efficiency are also observed different storytelling genres \cite{segel2010narrative}, such as data videos \cite{amini2015datavideos, xian2023datavideos} or data comics \cite{bach2017emerging, milesi2024comics}.

This paper contributes to the InfoVis and HCI literature by building on prior frameworks, such as those proposed by \citet{adar2021communicative} and \citet{shao2024datastorytelling}, to systematically assess the role of DS in supporting cognitive tasks across the full spectrum of complexity, as classified Bloom's taxonomy. Building on this theoretical foundation, we demonstrate its applicability to real-world tasks requiring both lower-order and higher-order skills through a large-scale empirical evaluation of DS-enhanced visualisations. Our study spans all levels of Bloom’s taxonomy, and provides insights into how DS can support higher-order cognitive levels -- \textit{Apply} to \textit{Create} -- which have been underexplored in the literature but are critical for communicating complex ideas in practice \cite{morth2023scrollyvis, seyser2018scrollytelling}. Unlike the work of \citet{leerobbins2022learning}, which focuses primarily on design specifications and user preferences, this study evaluates how these design principles translate into measurable improvements in task performance across varying levels of cognitive complexity. Additionally, we introduce a human-in-the-loop-AI task generation process, which combines large language models, domain-specific validation tools, and human oversight to create tasks aligned with Bloom's taxonomy. This approach provides a framework for designing cognitively diverse tasks, and similar to \citet{yan2024data,yan2024vizchat}, highlights the potential for human-AI collaboration in the HCI and InfoVis community.

\subsection{Implications for Design and Practice}
The different effects of DS-enhanced and conventional visualisations across various cognitive levels present distinct challenges and opportunities for researchers, practitioners, and designers. DS elements can significantly improve the effectiveness \textbf{(RQ1)} in understanding tasks and improve the efficiency \textbf{(RQ2)} in higher-order cognitive tasks requiring interpretation and critical thinking. Educators and trainers could consider adding DS elements based on specific learning objectives and the expected cognitive engagement levels of their audience. This could lead to higher learning efficiency and effectiveness, particularly in data-driven disciplines.

Despite the advantages of DS-enhanced visualisations in improving effectiveness in the \textit{Understanding} tasks and efficiency across most cognitive levels, it is noteworthy that in the \textit{Create} task, where users composed text based on the visualisations, participants using DS-enhanced produced fewer \textit{Evaluate} sentences (\textbf{RQ3}). Although this difference was not statistically significant, it warrants further investigation. This finding suggests a potential risk of DS-enhanced visualisations in contexts such as news or sensitive topics, where presenting a single narrative may inadvertently discourage critical thinking and reduce users' inclination to question underlying assumptions or consider missing information not presented in the graph. Understanding this effect more deeply could highlight an important limitation of data storytelling, especially when fostering critical evaluation is essential \cite{srinivasan2019augmenting}. This is consistent with the concerns raised by \citet{milesi2024storytelling}, who discussed the potential for critical data points or narratives to be inadvertently obscured by the inclusion of DS enhancements.

Based on the results of this study, we recommend the following:
\begin{itemize}
    \item For complex cognitive tasks, DS-enhancements did not support the correctness (\textbf{RQ1}) but tended to facilitate more \textit{efficient} information retrieval (\textbf{RQ2}). Consequently, conventional visualisations should be enhanced with DS elements in scenarios where the primary objective is to complete complex tasks more efficiently. 
    \item The relatively low level of critical evaluation in the DS-enhanced condition compared to the conventional visualisations (\textbf{RQ3}) suggests that while DS-enhancements can be leveraged to guide attention, they may lead to individuals interpreting the data more passively compared to those who explore the visualisations and discover the insights themselves. Therefore, depending on the context, it is important to balance the implementation of DS-enhancements to ensure that the intended message is still clear while also permitting critical evaluation of the data. For example, by using explicit elicitation techniques like contrastive narratives \cite{rogha2024impact} or interactive elements \cite{morth2023scrollyvis,shi2023breaking}. 
    \item Given that not all participants preferred DS-enhanced visualisations over conventional visualisations (\textbf{RQ4}), designers should consider creating adaptable visualisation tools that can modify the level of narrative complexity and data storytelling elements based on user preferences \cite{liu2020survey}, feedback, or the intended cognitive load of the task \cite{carenini2014highlighting}. As noted by \citet{figueiras2014narrative} if a user is proficient in the task, DS-enhanced visualisations may be perceived as ``\textit{boring}''. Conversely, users who are not proficient in the task may find highly exploratory visualisations difficult to engage with or comprehend \cite{figueiras2014narrative}. Therefore, the integration of user control features, such as the ability to toggle between more and less complex narrative features, may help accommodate a wider range of users and use cases.
\end{itemize}

\subsection{Limitations and Future Works}
A limitation of this study is the focus on line charts and maps as visualisation types. This choice was made to minimise potential fatigue due to the survey's length \cite{fereshteh2018evaluating} and to maintain consistency with the real-world context of \textit{Our World in Data}, whose preexisting data-driven narratives constrained the design to align with the information and presentation of the original visualisations. As a result, factors such as variations in contextual annotations (explored in \citet{stokes2023striking}) or DS-enhancements were not comprehensively explored. Instead, this study prioritised observing differences in task-solving performance across varying complexities within a defined DS-enhancement framework. Future studies should expand on these findings by exploring a broader range of visualisations (e.g., bar charts, scatter plots) and storytelling components. This study validates findings from studies that used synthetic data (e.g., \citet{stokes2023striking}) with real-world visualisations, serving as a foundation for future research to explore the generalisability of these results in more realistic contexts, such as full articles \cite{rogha2024impact} or interactive narratives \cite{shi2023breaking,morth2023scrollyvis}. For example, including breaking the fourth wall techniques, such as the \textit{Golden Hook} for capturing interest and the \textit{Magic Mirror} for tailoring insights as done in \citet{shi2023breaking}; as well as visual narrative flow elements including story layout and progression \cite{mckenna2017visual}, or sequencing techniques \cite{hullman2013deeper}.

While participants were instructed to not use external resources to answer questions and ``\textit{I am not sure}'' was given as an option for the multiple-choice questions, the online nature of this study means that it is possible that participants used search engines or GenAI to answer the questions. Despite efforts to exclude participants suspected of using GenAI in the written tasks, based on indicators such as words per minute, specific vocabulary, and grammatical structures, it remains difficult to completely verify the use of external sources. Consequently, conducting this study in a controlled in-person setting could mitigate the risk of cheating and obtain a richer qualitative dataset as participants work through the questions.

Another limitation of the study is the open-ended nature of the \textit{Create} task. No specific instructions were provided to participants outside of the general prompt asking them to write about the relationship between the visualisations. This left both the interpretation of the question and the content of their response to the discretion of the participants. As there is variation in the participants' demographics (age, background, country of origin, etc.) this may have influenced the responses to RQ3.
